\section{Scalar Edge and Entropy Inequalities}\label{app:scalar}\label{ck-the-scalar-and-edge-lemmas}

This appendix proves the scalar and edge facts used in
Section~\ref{sec:mechanism}. Write
\begin{align*}
 L&=\log2,& h(r)&=L-\phi(r),\\
 g(r)&=\operatorname{atanh}r,& F(v)&=\frac{h(\tanh v)}{\tanh v}.
\end{align*}
All these facts are independent of dimension.

\subsection{Convex Lower Bound for Centered Edges}\label{centered-edges-give-a-convex-lower-bound}

For edge values \(p,q\in(-1,1)\), define
\begin{align*}
 \edgewidth&=\frac{|p-q|}{2},& H_e&=\frac{h(p)+h(q)}2,\\
 a(p,q)&=\frac{(p-q)(g(p)-g(q))}{4},\\
 j(p,q)&=\frac{(\arcsin p-\arcsin q)^2}{4}.
\end{align*}
The edge contributes \(a(p,q)\) to the action and \(j(p,q)\) to the angle energy.
We abbreviate these edge functions as \(a,j\) within this subsection. The subscript in \(\edgewidth\) names its role: edge half-width.
Let \(F^{-1}\) denote the inverse function of \(F\).
For a positive scalar \(k\), first define
\[
 \theta_{\rm edge}(k)=\arcsin(\tanh k),\qquad
 \corr(F(k))=k-\frac{\theta_{\rm edge}(k)^2}{\tanh k}.
\]
Here \(\corr\) is short for the retained edge correction.
We will prove
\begin{align}
 a(p,q)&\ge \edgewidth F^{-1}(H_e/\edgewidth),\nonumber\\
 a(p,q)-j(p,q)&\ge \edgewidth\corr(H_e/\edgewidth).
 \tag{S2}\label{eq:scalar-S2}
\end{align}
The right-hand sides are convex functions of \((\edgewidth,H_e)\) and
decrease with \(H_e\). This permits averaging over all edges in a
coordinate. At \(\edgewidth=0\), both expressions are defined by their
limiting value zero.

\subsubsection{Centering property}\label{why-centering-the-edge-helps}

For any real number \(r\), define
\(\sinh r=(e^r-e^{-r})/2\) and \(\cosh r=(e^r+e^{-r})/2\).
For \(r\ne0\), \(\coth r=\cosh r/\sinh r\).
In this edge calculation \(v\) is a real center and \(k>0\) a
half-width, specified by \(p=\tanh(v+k)\), \(q=\tanh(v-k)\).
Assume \(p>q\). Then \(a=\edgewidth k\). Set \(M=\cosh(2v)\) and
\(C_k=\cosh(2k)\). Direct calculation gives
\begin{align*}
 \edgewidth&=\frac{\sinh(2k)}{M+C_k},\\
 H_e&=\frac12\log[2(M+C_k)]
       -\frac{v\sinh(2v)+k\sinh(2k)}{M+C_k}.
\end{align*}
To differentiate with respect to the center, first consider \(v>0\)
and set \(x=2v\). The derivative of the numerator of \(H_e/\edgewidth \) has the sign of \[
 \log[2(\cosh x+C_k)]-x\coth x.
\]
It is nonnegative because, for \(C_k\ge1\),
\begin{align*}
 \log[2(\cosh x+C_k)]
 &\ge x+2\log(1+e^{-x})\\
 &\ge x+\frac2{e^x+1}\\
 &\ge x+\frac{2x}{e^{2x}-1}=x\coth x.
\end{align*}
The last two steps use \(\log(1+y)\ge y/(1+y)\) for \(y\ge0\)
and \(e^x-1\ge x\). Evenness handles \(v<0\). Hence \[
                         H_e/\edgewidth \ge F(k).               \tag{S3}\label{eq:scalar-S3}
\]

The function \(\arctan x\) is the angle between \(-\pi/2\) and
\(\pi/2\) whose tangent is \(x\).
For the angle term, let \(b_{\rm edge}=\sinh k\) and \(x=b_{\rm edge}/\cosh v\).
This abbreviation is used only in the following half-angle calculation.
The half-angle identity gives
\begin{align*}
 \arcsin p-\arcsin q&=2\arctan x,\\
 \frac{j}{\edgewidth}
 &=\frac{b_{\rm edge}}{\sqrt{1+b_{\rm edge}^2}}(1+x^2)
   \left(\frac{\arctan x}{x}\right)^2.
\end{align*}
The last factor increases with \(x\): its logarithmic derivative is
\(2(x-\arctan x)/[x(1+x^2)\arctan x]>0\). Since
\(x\le b_{\rm edge}\), it follows that
\[
                         j/\edgewidth \le\theta_{\rm edge}(k)^2/\tanh k.    \tag{S4}\label{eq:scalar-S4}
\]

\subsubsection{Averaging by convexity}\label{why-the-bounds-can-be-averaged}

Put \(\ell=\log(2\cosh k)\). Differentiation gives \[
 F'(k)=-\frac{\ell}{\sinh^2 k},\qquad
 (F^{-1})'(F(k))=-\frac{\sinh^2 k}{\ell}.
\] The positive ratio on the right increases with \(k\), since \(2\ell>\tanh^2 k\). Thus \(F^{-1}\) is decreasing and convex.

Similarly,
\begin{align*}
 \frac d{dk}\left(k-\frac{\theta_{\rm edge}(k)^2}{\tanh k}\right)
 &=\left(1-\frac{\theta_{\rm edge}(k)}{\sinh k}\right)^2,\\
 \corr'(F(k))&=-\frac{(\sinh k-\theta_{\rm edge}(k))^2}{\ell}.
\end{align*}
The positive ratio in the last expression increases with \(k\): its
derivative has the sign of
\(2\ell\sinh k-(\sinh k-\theta_{\rm edge}(k))>0\). Therefore \(\corr\)
is nonnegative, decreasing, and convex.

Combining these facts with \eqref{eq:scalar-S3}--\eqref{eq:scalar-S4}
proves \eqref{eq:scalar-S2}.
Let \(\psi\) be either of the two functions \(F^{-1}\) or \(\corr\).
The function of two variables \(\edgewidth \psi(H_e/\edgewidth )\), called its
\emph{perspective}, is convex, increases with \(\edgewidth \), and decreases
with \(H_e\). Jensen may therefore be applied to \eqref{eq:scalar-S2}.

For a monotone field \(u=T_\rho f\), a coordinate \(i\) with \(a_i>0\) satisfies
\(\E\edgewidth=\rho a_i\) and \(\E H_e=\E h(u)\).
For any upper bound \(H_{\rm budget}\ge\E h(u)>0\), Jensen and
decrease in entropy give
\begin{align*}
 \E a&\ge\rho a_i v_i,\\
 \E(a-j)&\ge\rho a_i[v_i-c(v_i)],\\
 F(v_i)&=H_{\rm budget}/(\rho a_i).
\end{align*}
Summing proves Lemma~\ref{lem:edges}. Its two applications simply
supply different entropy budgets.

\subsection{Monotonicity of the Entropy Profile}
\label{a-short-proof-that-vfv-decreases}

\begin{appendixlemma}[Profile monotonicity]
\label{lem:global-profile}
The function \(vF(v)\) strictly decreases for \(v>0\).
\end{appendixlemma}

One concavity argument supplies the needed entropy bound. The function
\((1-r^2)g(r)-r h(r)\) has continuous value zero at both endpoints
of \([0,1]\), and its second derivative is \(-r/(1-r^2)<0\).
It is therefore positive inside the interval, giving
\[
                   h(r)<\frac{(1-r^2)g(r)}r.
\]
Now fix \(v>0\) and put \(r=\tanh v\), so \(g(r)=v\).
The factor \(r-v(1-r^2)\) is positive: it vanishes at zero
and its derivative in \(v\) is \(2vr(1-r^2)>0\).
We may consequently use the entropy bound in the derivative
\begin{align*}
 \frac{d}{dv}[vF(v)]
 &=\frac{h(r)[r-v(1-r^2)]}{r^2}-\frac{v^2(1-r^2)}r\\
 &<\frac{v(1-r^2)}{r^3}(r-v)<0,
\end{align*}
where the last step uses \(\tanh v<v\).

\paragraph{Consequence for the small-height comparison.}
Choose an upper height \(0<\cutoff\le\logodds\). For \(0<v\le\cutoff\) and
\(a=F(\logodds)/F(v)\), monotonicity gives
\(a/v\le F(\logodds)/[\cutoff F(\cutoff)]\).
Also \(c(v)\le v\), since \(\corr\ge0\).
The LSI branch of Lemma~\ref{lem:deficit} now gives
\[
 \unitcost(v)\le\rho+p\Klsi\rho^2\frac{F(\logodds )}{\cutoff F(\cutoff)}
                   \qquad(0<v\le\cutoff).                 \tag{S5}\label{eq:scalar-S5}
\] This bounds the entire small-\(v\) interval by one expression.

\subsection{Relative Rates of the Entropy Profile and Angle Correction}
\label{app:profile-angle-rates}

The middle argument compares the profile \(vF(v)\) with the edge
correction \(v-c(v)\). The following statement involves only these
two functions, independently of the channel or the Boolean rule.

\begin{appendixlemma}[The profile pays for the angle correction]
\label{lem:profile-angle-rates}
Define the logarithmic decay rate \(Q(v)=-F'(v)/F(v)\). Then
\begin{equation}
 \begin{aligned}
 Q(v)-\frac1v&>\frac32[1-c'(v)]&& (v>0),\\
 Q(v)-\frac1v&>2[1-c'(v)]&& (0<v\le2).
 \end{aligned}
 \label{eq:middle-profile-angle-rates}
\end{equation}
Equivalently, the rate of decrease of \(\log[vF(v)]\) exceeds
three halves of the rate of increase of \(v-c(v)\); up to height
two, the factor improves to two.
\end{appendixlemma}

\begin{proof}
We first record the angle geometry used in the rate comparison.
For a variable height \(v>0\), write
\(r_v=\tanh v\), \(\tau_v=\sech v\),
\(\theta_v=\arcsin r_v\), and \(s(v)=\theta_v/\sinh v\).
Since \(\theta_v'=\tau_v\), we have
\(0<\theta_v<v<\sinh v\), hence \(0<s<1\). Differentiation gives
\[
 s'=\frac{\tau_v^2-s}{r_v}<0,\qquad
 c'=2s-s^2,\qquad c''=2(1-s)s'<0.
\]
The first sign follows from \(\theta_v>r_v\tau_v\).
Thus \(c\) increases, and \(c(v)/v\) decreases because \(c\) is
concave and \(c(0)=0\).

We will use the two angle bounds
\begin{equation}
 \frac{3z}{2+\sqrt{1-z^2}}\le\arcsin z
 \le\frac{\pi z}{2+\sqrt{1-z^2}}
 \qquad(0\le z\le1).
 \label{eq:middle-angle-brackets}
\end{equation}
Here is a proof of both at once. Set \(z=\sin t\),
\(0<t<\pi/2\). The derivative of
\(t(2+\cos t)/\sin t\) has numerator
\((2+\cos t)\sin t-t(1+2\cos t)\). This numerator vanishes at
zero and has derivative \(2\sin t(t-\sin t)>0\). The quotient
therefore increases from three to \(\pi\), proving the bounds.

Since \(1-c'=(1-s)^2>0\), it suffices to compare
\(Q-1/v\) with the indicated multiples of \((1-s)^2\).
The fixed estimate \(e^4<55\) below follows from the finite sums in
Section~\ref{app:fixed-arithmetic}.

Two simple lower bounds on \(Q\) suffice. With \(q_v=e^{-2v}\),
the entropy formula gives
\[
 Q(v)=\frac{4q_v[v+\log(1+q_v)]}
 {(1-q_v)[(1+q_v)\log(1+q_v)+2vq_v]}
 \ge\frac{4v}{2v+1}.
\]
Indeed \(\log(1+q_v)\le q_v\) bounds the denominator by
\(q_v(1+2v)\), while the numerator is at least \(4q_vv\).
Also \(Q(v)>\coth v\). To see this, let
\begin{align*}
 \mathcal H(v)&=\cosh v\,h(\tanh v),\\
 D(v)&=v\coth v-\log(2\cosh v).
\end{align*}
We have \(D'=(\tanh v-v)/\sinh^2v<0\) and
\(D(v)\longrightarrow0\) at infinity, so
\(\mathcal H'=-\sinh v\,D<0\). Since
\(F=\mathcal H/\sinh v\), this proves the claimed bound.

The positive power series show that \((\cosh v-1)/v^2\) and
\(\sinh v/v\) increase. The geometric estimates
\begin{align*}
 \cosh1-1&<\frac{1/2}{1-1/12}<\frac35,\\
 \sinh1-1&<\frac{1/6}{1-1/20}<\frac15,\\
 \cosh2-1&<\frac2{1-1/3}=3
\end{align*}
and \(\cosh4<(55+1)/2<33\) will give convenient quadratic bounds.
For \(0<v\le1\), they imply
\(\cosh v\le1+3v^2/5\). By
\eqref{eq:middle-angle-brackets},
\[
 1-s\le\frac{2v^2}{5+2v^2},\qquad
 Q(v)-1/v>\coth v-1/v>\frac{5v}{18}>\frac v4.
\]
The second inequality uses
\(v\cosh v-\sinh v=\int_0^v t\sinh t\,dt\ge v^3/3\) and
\(\sinh v<6v/5\). Finally
\((5+2v^2)^2-32v^3\ge25-12v^2>0\), so
\(v/4>2[2v^2/(5+2v^2)]^2\).

For \(1\le v\le2\), use \(\cosh v\le1+3v^2/4\), giving
\(1-s\le v^2/(2+v^2)\). The numerator of
\[
 \frac{4v}{2v+1}-\frac1v-2\frac{v^4}{(2+v^2)^2},
\]
after multiplication by \(v(2v+1)(2+v^2)^2\), is
\[
 v^3(15v-4v^2-8)+4(3v+1)(v-1)>0.
\]
The concave quadratic in the first term is at least three on
\([1,2]\); the second term is nonnegative.
This proves the stronger rate bound for \(0<v\le2\).

For \(2\le v\le4\), the bound \(\cosh v\le1+2v^2\) gives
\(1-s\le4v^2/(3+4v^2)\). The numerator of
\[
 \frac{4v}{2v+1}-\frac1v-\frac32
                   \frac{16v^4}{(3+4v^2)^2},
\]
after multiplication by \(v(2v+1)(3+4v^2)^2\), is
\[
 16v^4(v-2)(v-3/2)+16v^3(2v-3)+3(4v^2-6v-3)>0.
\]
Each term is nonnegative, and the last quadratic is increasing
from one. For \(v\ge4\), simply use
\(Q(v)-1/v\ge2-2/(2v+1)-1/v>3/2\) and \((1-s)^2<1\).
This proves the global rate bound.

\end{proof}

To use this lemma, integrate its rate bound between a coordinate
height and the channel height. Appendix~\ref{app:middle-channel}
checks which coefficient the channel needs; the profile calculation
itself is now complete.


\subsection{Tensorized Logarithmic-Sobolev Inequality}
\label{app:lsi-proof}

The coordinate estimate in Lemma~\ref{lem:deficit} uses
\[
             \Ent(w^2)\le2\E[w\Lapl w].
\]
This classical inequality measures how far the nonnegative field
\(w^2\) is from constant by the squared differences of \(w\) across
edges. We include its proof, including the step that makes its constant
independent of dimension.

Start with one bit. Replacing the two values of \(w\) by their absolute
values preserves entropy and reduces their squared difference.
Homogeneity then lets us assume \(\E w^2=1\) and write the values as
\(\sqrt{1+s},\sqrt{1-s}\), where \(0\le s\le1\). In this notation
\begin{align*}
 \Ent(w^2)&=\phi(s),\\
 2\E[w\Lapl w]&=\frac{(w_+-w_-)^2}{2}=1-\sqrt{1-s^2}.
\end{align*}
The derivative of \(s/\sqrt{1-s^2}-\atanh s\) is
\((1-s^2)^{-3/2}-(1-s^2)^{-1}\ge0\). The difference vanishes at
zero. Integrating \(\phi'(s)=\atanh s\) therefore proves
\(\phi(s)\le1-\sqrt{1-s^2}\). Continuity supplies \(s=1\).

It remains to tensorize the one-bit bounds. For a finite probability
space, entropy is convex as a function of a nonnegative vector \(z\).
For positive \(z\), its second derivative in a direction \(v\) is
\[
 \E\frac{v^2}{z}-\frac{(\E v)^2}{\E z}\ge0
\]
by Cauchy--Schwarz. Zero entries follow by continuity.
For independent coordinates \(X,Y\), the entropy chain rule and this
convexity give
\begin{align*}
 \Ent_{X,Y}(z)
 &=\E_Y\Ent_X(z)+\Ent_Y(\E_Xz)\\
 &\le\E_Y\Ent_X(z)+\E_X\Ent_Y(z).
\end{align*}
Induction over the coordinates yields
\(\Ent(z)\le\sum_i\E\Ent_i(z)\), where \(\Ent_i\) takes entropy
over coordinate \(i\) with all others fixed. Apply the one-bit result
on each such pair, with \(z=w^2\):
\[
 \Ent(w^2)\le\frac12\sum_i\E_{-i}(w_{i,+}-w_{i,-})^2
                         =2\E[w\Lapl w].
\]
No constant was multiplied by the number of coordinates.

\subsection{Propagation of the Local Cutoff}
\label{app:local-propagation}
This is the threshold-propagation property in
Yu's Lemma 4.9~\cite{Yu}, expressed in sign means and natural logarithms.
We include the one-sign-change proof so that the local comparison
requires no numerical monotonicity test or external lemma.

Define the scalar local margin
\[
 M_{\rm loc}(r,a)=(1-r^2)h(ar)-(1-ar^2)h(r),\qquad 0\le r,a\le1.
\]
We prove the propagation used below: if \(M_{\rm loc}(r_1,a)\ge0\) for some
\(0<r_1<1\), then \(M_{\rm loc}(r,a)\ge0\) whenever \(0\le r\le r_1\).
The reason is a single change of sign in its entropy-series coefficients.

The case \(a=1\) is an identity. For \(a<1\), put \(t=r^2\) and
\(H(t)=h(\sqrt t)\). Let \(b_n=1/[2n(2n-1)]\) be the entropy-series
coefficient, and let \(T_n=\sum_{j>n}b_j\) be its tail. Then
\begin{align*}
 H(t)&=L-\sum_{n\ge1}b_nt^n,\\
 \frac{H(t)}{1-t}&=\sum_{n\ge0}T_nt^n,\\
 T_n&=\int_0^1\frac{x^{2n}}{1+x}\,\mathrm dx.
\end{align*}
For the last identity, express
\(b_j=\int_0^1x^{2j-2}(1-x)\,\mathrm dx\), sum the geometric series,
and cancel \(1-x\). All series used here converge absolutely for
\(t<1\).

Expansion of the local margin now gives
\begin{align*}
 \frac{M_{\rm loc}(\sqrt t,a)}{1-t}&=\sum_{n\ge1}q_nt^n,\\
 q_n&=\frac{1-a}{2n}\left[
 \frac{a+a^2+\cdots+a^{2n-1}}{2n-1}-2nT_n\right].
\end{align*}
The average of the powers of \(a\) decreases with \(n\): the next
average appends two terms no larger than the preceding terms.
In the opposite direction,
\[
 (n+1)T_{n+1}-nT_n
 =T_n-\frac1{2(2n+1)}>0,
\]
as is seen by replacing \(1/(1+x)\) by \(1/2\) in the integral.
Thus the coefficients \(q_n\) are nonnegative first and negative
thereafter, with at most one change of sign. They eventually are
negative, since the power average tends to zero whereas
\(2nT_n>n/(2n+1)\).

Let \(N\) be the last index with \(q_N\ge0\). Such an index exists
if the margin is nonnegative at \(r_1>0\). For
\(t_1=r_1^2\) and \(\lambda=t/t_1\in(0,1]\), the sign ordering gives
\[
 \sum_{n\ge1}q_nt^n\ge
             \lambda^N\sum_{n\ge1}q_nt_1^n\ge0.
\]
Indeed, \(\lambda^n\ge\lambda^N\) before the sign change; afterwards
the inequality reverses when multiplied by the negative coefficient.
This proves the propagation, with \(r=0\) immediate.

Finally \(\partial_a^2M_{\rm loc}(r,a)=-(1-r^2)r^2/(1-a^2r^2)<0\) and
\(M_{\rm loc}(r,1)=0\). A successful check at \((r_1,a_0)\) therefore covers
all \(r\le r_1\) and \(a\ge a_0\): concavity extends in \(a\),
and the series argument extends in \(r\).


\subsection{The Influential-Coordinate Criterion}\label{a-large-singleton-coefficient-already-proves-ck}

The starting point is the standard strong data-processing bound
obtained from cube logarithmic Sobolev~\cite{Gross,ODonnell}.
The local-optimality perspective is shared with~\cite{Sam,Yu}; the
calculation below makes explicit the mean correction needed to use
this particular criterion for biased rules.

For an arbitrary Boolean \(f\), fix a coordinate and write \(a=|\widehat f(\{i\})|\), \(m=Ef\), \(\delta=1-\rho^2\). Necessarily \(|m|+a\le1\).

For a cube field \(v\) taking values in \([-1,1]\), define
\(\operatorname{Ent}_{\phi}(v)=\E\phi(v)-\phi(\E v)\).
We first derive the entropy contraction used here and in Appendix~\ref{app:low}:
\[
 \operatorname{Ent}_{\phi}(T_\rho v)
                 \le\rho^2\operatorname{Ent}_{\phi}(v).       \tag{S6a}\label{eq:scalar-S6a}
\]
For a positive field \(w\), the elementary edge inequality
\[
 (p-q)(\log p-\log q)\ge4(\sqrt p-\sqrt q)^2
 \qquad(p,q>0)
\]
follows by Cauchy--Schwarz applied to the integrals of \(1\) and
\(1/\sqrt t\) between \(p\) and \(q\). Sum over cube edges and use the
logarithmic Sobolev inequality just proved in
Section~\ref{app:lsi-proof}:
\[
 \E[(\Lapl w)\log w]\ge4\E[\sqrt w\,\Lapl \sqrt w]\ge2\Ent(w).
\]
Let \(t\ge0\) be the noise-flow time and define \(w_t=T_{e^{-t}}w\).
The mean is constant and
\(\frac{\diff}{\diff t}\Ent(w_t)=-\E[(\Lapl w_t)\log w_t]\).
Integration gives
\(\Ent(T_\rho w)\le\rho^2\Ent(w)\).
For a positive real \(\varepsilon\), nonnegative fields follow
by replacing \(w\) with \(w+\varepsilon\)
and taking \(\varepsilon\downarrow0\).
Finally,
\[
 \operatorname{Ent}_{\phi}(v)
       =\tfrac12\{\Ent(1+v)+\Ent(1-v)\},
\]
which proves \eqref{eq:scalar-S6a}.

Apply \eqref{eq:scalar-S6a} to the other coordinates after noising the
chosen bit first.
The intermediate field has values \(\pm\rho\) on pivotal edges and
\(\pm1\) otherwise. Pivotal probability is at least \(a\), and the section
means are \(m\pm\rho a\). Thus \[
 Eh(T_\rho f)\ge\rho^2a\,h(\rho)
                      +\frac{\delta}{2}[h(m+\rho a)+h(m-\rho a)]. \tag{S6}\label{eq:scalar-S6}
\]

The mean costs no additional loss because \[
 \phi(m)+\frac{\delta}{2}[h(m+\rho a)+h(m-\rho a)]
                              \ge\delta h(\rho a).    \tag{S7}\label{eq:scalar-S7}
\] For completeness, twice differentiate the difference in \(m\). After positive factors are removed, its second derivative has the sign of \[
 [1-(m+a)^2][1-(m-a)^2]
          +(1-\rho^2)a^2(1-m^2-a^2)\ge0.
\] The difference and its first derivative vanish at zero. Integrating
proves \eqref{eq:scalar-S7}; reflection handles negative \(m\).

Combining \eqref{eq:scalar-S6}--\eqref{eq:scalar-S7} proves CK whenever \[
                 (1-\rho^2)h(\rho a)
                    -(1-\rho^2a)h(\rho)\ge0.          \tag{S8}\label{eq:scalar-S8}
\] For fixed \(\rho\), the left side is concave in \(a\) and zero at
\(a=1\). Checking it at a cutoff therefore covers every larger \(a\).
Section~\ref{app:local-propagation} proves directly why a successful
cutoff also works at smaller correlations. The all-bias step is
\eqref{eq:scalar-S7}, proved here.
